\section{}

\paragraph{1127年}\eventref{SS-1127-EVENT-111}{SS-1127-EVENT-111}　赵构即位为高宗，南宋朝廷开始重建。此项以《宋史》卷24—47〈本纪〉；《建炎以来系年要录》本年条；《宋会要辑稿》相关门类 为基本来源线索；编年与官修材料能够确认年序和制度用语，但对地方执行、人数、动机或社会后果的判断仍须同文书、碑刻、考古及对方叙事互校。

\paragraph{1127年}\eventanchor{SS-1127-REGNAL-YEAR}{SS-1127-REGNAL-YEAR}　南宋以宋高宗在位第1年作为诏令、赋役与外交文书的纪年框架。此项以《宋史》卷24—47〈本纪〉；《建炎以来系年要录》本年条；《宋会要辑稿》相关门类 为基本来源线索；编年与官修材料能够确认年序和制度用语，但对地方执行、人数、动机或社会后果的判断仍须同文书、碑刻、考古及对方叙事互校。

\paragraph{1128年}\eventanchor{SS-1128-EVENT-112}{SS-1128-EVENT-112}　高宗朝廷驻跸扬州等地，组织勤王与财政。此项以《宋史》卷24—47〈本纪〉；《建炎以来系年要录》本年条；《宋会要辑稿》相关门类 为基本来源线索；编年与官修材料能够确认年序和制度用语，但对地方执行、人数、动机或社会后果的判断仍须同文书、碑刻、考古及对方叙事互校。

\paragraph{1128年}\eventanchor{SS-1128-REGNAL-YEAR}{SS-1128-REGNAL-YEAR}　南宋以宋高宗在位第2年作为诏令、赋役与外交文书的纪年框架。此项以《宋史》卷24—47〈本纪〉；《建炎以来系年要录》本年条；《宋会要辑稿》相关门类 为基本来源线索；编年与官修材料能够确认年序和制度用语，但对地方执行、人数、动机或社会后果的判断仍须同文书、碑刻、考古及对方叙事互校。

\paragraph{1129年}\eventanchor{SS-1129-EVENT-113}{SS-1129-EVENT-113}　苗傅、刘正彦发动兵变，随后被平定。此项以《宋史》卷24—47〈本纪〉；《建炎以来系年要录》本年条；《宋会要辑稿》相关门类 为基本来源线索；编年与官修材料能够确认年序和制度用语，但对地方执行、人数、动机或社会后果的判断仍须同文书、碑刻、考古及对方叙事互校。

\paragraph{1129年}\eventanchor{SS-1129-REGNAL-YEAR}{SS-1129-REGNAL-YEAR}　南宋以宋高宗在位第3年作为诏令、赋役与外交文书的纪年框架。此项以《宋史》卷24—47〈本纪〉；《建炎以来系年要录》本年条；《宋会要辑稿》相关门类 为基本来源线索；编年与官修材料能够确认年序和制度用语，但对地方执行、人数、动机或社会后果的判断仍须同文书、碑刻、考古及对方叙事互校。

\paragraph{1130年}\eventanchor{SS-1130-EVENT-114}{SS-1130-EVENT-114}　金军渡江后北撤，南宋保住江南核心地区。此项以《宋史》卷24—47〈本纪〉；《建炎以来系年要录》本年条；《宋会要辑稿》相关门类 为基本来源线索；编年与官修材料能够确认年序和制度用语，但对地方执行、人数、动机或社会后果的判断仍须同文书、碑刻、考古及对方叙事互校。

\paragraph{1130年}\eventanchor{SS-1130-REGNAL-YEAR}{SS-1130-REGNAL-YEAR}　南宋以宋高宗在位第4年作为诏令、赋役与外交文书的纪年框架。此项以《宋史》卷24—47〈本纪〉；《建炎以来系年要录》本年条；《宋会要辑稿》相关门类 为基本来源线索；编年与官修材料能够确认年序和制度用语，但对地方执行、人数、动机或社会后果的判断仍须同文书、碑刻、考古及对方叙事互校。

\paragraph{1131年}\eventanchor{SS-1131-REGNAL-YEAR}{SS-1131-REGNAL-YEAR}　南宋以宋高宗在位第5年作为诏令、赋役与外交文书的纪年框架。此项以《宋史》卷24—47〈本纪〉；《建炎以来系年要录》本年条；《宋会要辑稿》相关门类 为基本来源线索；编年与官修材料能够确认年序和制度用语，但对地方执行、人数、动机或社会后果的判断仍须同文书、碑刻、考古及对方叙事互校。

\paragraph{1132年}\eventanchor{SS-1132-EVENT-115}{SS-1132-EVENT-115}　行在逐步以临安为稳定基地。此项以《宋史》卷24—47〈本纪〉；《建炎以来系年要录》本年条；《宋会要辑稿》相关门类 为基本来源线索；编年与官修材料能够确认年序和制度用语，但对地方执行、人数、动机或社会后果的判断仍须同文书、碑刻、考古及对方叙事互校。

\paragraph{1132年}\eventanchor{SS-1132-REGNAL-YEAR}{SS-1132-REGNAL-YEAR}　南宋以宋高宗在位第6年作为诏令、赋役与外交文书的纪年框架。此项以《宋史》卷24—47〈本纪〉；《建炎以来系年要录》本年条；《宋会要辑稿》相关门类 为基本来源线索；编年与官修材料能够确认年序和制度用语，但对地方执行、人数、动机或社会后果的判断仍须同文书、碑刻、考古及对方叙事互校。

\paragraph{1133年}\eventanchor{SS-1133-REGNAL-YEAR}{SS-1133-REGNAL-YEAR}　南宋以宋高宗在位第7年作为诏令、赋役与外交文书的纪年框架。此项以《宋史》卷24—47〈本纪〉；《建炎以来系年要录》本年条；《宋会要辑稿》相关门类 为基本来源线索；编年与官修材料能够确认年序和制度用语，但对地方执行、人数、动机或社会后果的判断仍须同文书、碑刻、考古及对方叙事互校。

\paragraph{1134年}\eventanchor{SS-1134-EVENT-116}{SS-1134-EVENT-116}　南宋军在襄汉、淮西等方向取得局部胜利。此项以《宋史》卷24—47〈本纪〉；《建炎以来系年要录》本年条；《宋会要辑稿》相关门类 为基本来源线索；编年与官修材料能够确认年序和制度用语，但对地方执行、人数、动机或社会后果的判断仍须同文书、碑刻、考古及对方叙事互校。

\paragraph{1134年}\eventanchor{SS-1134-REGNAL-YEAR}{SS-1134-REGNAL-YEAR}　南宋以宋高宗在位第8年作为诏令、赋役与外交文书的纪年框架。此项以《宋史》卷24—47〈本纪〉；《建炎以来系年要录》本年条；《宋会要辑稿》相关门类 为基本来源线索；编年与官修材料能够确认年序和制度用语，但对地方执行、人数、动机或社会后果的判断仍须同文书、碑刻、考古及对方叙事互校。

\paragraph{1135年}\eventanchor{SS-1135-EVENT-117}{SS-1135-EVENT-117}　岳飞等将领获得更大经略与军队整编责任。此项以《宋史》卷24—47〈本纪〉；《建炎以来系年要录》本年条；《宋会要辑稿》相关门类 为基本来源线索；编年与官修材料能够确认年序和制度用语，但对地方执行、人数、动机或社会后果的判断仍须同文书、碑刻、考古及对方叙事互校。

\paragraph{1135年}\eventanchor{SS-1135-REGNAL-YEAR}{SS-1135-REGNAL-YEAR}　南宋以宋高宗在位第9年作为诏令、赋役与外交文书的纪年框架。此项以《宋史》卷24—47〈本纪〉；《建炎以来系年要录》本年条；《宋会要辑稿》相关门类 为基本来源线索；编年与官修材料能够确认年序和制度用语，但对地方执行、人数、动机或社会后果的判断仍须同文书、碑刻、考古及对方叙事互校。

\paragraph{1136年}\eventanchor{SS-1136-REGNAL-YEAR}{SS-1136-REGNAL-YEAR}　南宋以宋高宗在位第10年作为诏令、赋役与外交文书的纪年框架。此项以《宋史》卷24—47〈本纪〉；《建炎以来系年要录》本年条；《宋会要辑稿》相关门类 为基本来源线索；编年与官修材料能够确认年序和制度用语，但对地方执行、人数、动机或社会后果的判断仍须同文书、碑刻、考古及对方叙事互校。

\paragraph{1137年}\eventanchor{SS-1137-REGNAL-YEAR}{SS-1137-REGNAL-YEAR}　南宋以宋高宗在位第11年作为诏令、赋役与外交文书的纪年框架。此项以《宋史》卷24—47〈本纪〉；《建炎以来系年要录》本年条；《宋会要辑稿》相关门类 为基本来源线索；编年与官修材料能够确认年序和制度用语，但对地方执行、人数、动机或社会后果的判断仍须同文书、碑刻、考古及对方叙事互校。

\paragraph{1138年}\eventanchor{SS-1138-EVENT-118}{SS-1138-EVENT-118}　临安被确立为行在和事实都城。此项以《宋史》卷24—47〈本纪〉；《建炎以来系年要录》本年条；《宋会要辑稿》相关门类 为基本来源线索；编年与官修材料能够确认年序和制度用语，但对地方执行、人数、动机或社会后果的判断仍须同文书、碑刻、考古及对方叙事互校。

\paragraph{1138年}\eventanchor{SS-1138-REGNAL-YEAR}{SS-1138-REGNAL-YEAR}　南宋以宋高宗在位第12年作为诏令、赋役与外交文书的纪年框架。此项以《宋史》卷24—47〈本纪〉；《建炎以来系年要录》本年条；《宋会要辑稿》相关门类 为基本来源线索；编年与官修材料能够确认年序和制度用语，但对地方执行、人数、动机或社会后果的判断仍须同文书、碑刻、考古及对方叙事互校。

\paragraph{1139年}\eventanchor{SS-1139-REGNAL-YEAR}{SS-1139-REGNAL-YEAR}　南宋以宋高宗在位第13年作为诏令、赋役与外交文书的纪年框架。此项以《宋史》卷24—47〈本纪〉；《建炎以来系年要录》本年条；《宋会要辑稿》相关门类 为基本来源线索；编年与官修材料能够确认年序和制度用语，但对地方执行、人数、动机或社会后果的判断仍须同文书、碑刻、考古及对方叙事互校。

\paragraph{1140年}\eventanchor{SS-1140-EVENT-119}{SS-1140-EVENT-119}　南宋军在河南、淮西等方向反攻。此项以《宋史》卷24—47〈本纪〉；《建炎以来系年要录》本年条；《宋会要辑稿》相关门类 为基本来源线索；编年与官修材料能够确认年序和制度用语，但对地方执行、人数、动机或社会后果的判断仍须同文书、碑刻、考古及对方叙事互校。

\paragraph{1140年}\eventanchor{SS-1140-REGNAL-YEAR}{SS-1140-REGNAL-YEAR}　南宋以宋高宗在位第14年作为诏令、赋役与外交文书的纪年框架。此项以《宋史》卷24—47〈本纪〉；《建炎以来系年要录》本年条；《宋会要辑稿》相关门类 为基本来源线索；编年与官修材料能够确认年序和制度用语，但对地方执行、人数、动机或社会后果的判断仍须同文书、碑刻、考古及对方叙事互校。

\paragraph{1141年}\eventanchor{SS-1141-EVENT-120}{SS-1141-EVENT-120}　南宋接受绍兴和议框架。此项以《宋史》卷24—47〈本纪〉；《建炎以来系年要录》本年条；《宋会要辑稿》相关门类 为基本来源线索；编年与官修材料能够确认年序和制度用语，但对地方执行、人数、动机或社会后果的判断仍须同文书、碑刻、考古及对方叙事互校。

\paragraph{1141年}\eventanchor{SS-1141-REGNAL-YEAR}{SS-1141-REGNAL-YEAR}　南宋以宋高宗在位第15年作为诏令、赋役与外交文书的纪年框架。此项以《宋史》卷24—47〈本纪〉；《建炎以来系年要录》本年条；《宋会要辑稿》相关门类 为基本来源线索；编年与官修材料能够确认年序和制度用语，但对地方执行、人数、动机或社会后果的判断仍须同文书、碑刻、考古及对方叙事互校。

\paragraph{1142年}\eventanchor{SS-1142-EVENT-121}{SS-1142-EVENT-121}　岳飞被处死，相关将领和军队受到整顿。此项以《宋史》卷24—47〈本纪〉；《建炎以来系年要录》本年条；《宋会要辑稿》相关门类 为基本来源线索；编年与官修材料能够确认年序和制度用语，但对地方执行、人数、动机或社会后果的判断仍须同文书、碑刻、考古及对方叙事互校。

\paragraph{1142年}\eventanchor{SS-1142-REGNAL-YEAR}{SS-1142-REGNAL-YEAR}　南宋以宋高宗在位第16年作为诏令、赋役与外交文书的纪年框架。此项以《宋史》卷24—47〈本纪〉；《建炎以来系年要录》本年条；《宋会要辑稿》相关门类 为基本来源线索；编年与官修材料能够确认年序和制度用语，但对地方执行、人数、动机或社会后果的判断仍须同文书、碑刻、考古及对方叙事互校。

\paragraph{1143年}\eventanchor{SS-1143-REGNAL-YEAR}{SS-1143-REGNAL-YEAR}　南宋以宋高宗在位第17年作为诏令、赋役与外交文书的纪年框架。此项以《宋史》卷24—47〈本纪〉；《建炎以来系年要录》本年条；《宋会要辑稿》相关门类 为基本来源线索；编年与官修材料能够确认年序和制度用语，但对地方执行、人数、动机或社会后果的判断仍须同文书、碑刻、考古及对方叙事互校。

\paragraph{1144年}\eventanchor{SS-1144-REGNAL-YEAR}{SS-1144-REGNAL-YEAR}　南宋以宋高宗在位第18年作为诏令、赋役与外交文书的纪年框架。此项以《宋史》卷24—47〈本纪〉；《建炎以来系年要录》本年条；《宋会要辑稿》相关门类 为基本来源线索；编年与官修材料能够确认年序和制度用语，但对地方执行、人数、动机或社会后果的判断仍须同文书、碑刻、考古及对方叙事互校。

\paragraph{1145年}\eventanchor{SS-1145-REGNAL-YEAR}{SS-1145-REGNAL-YEAR}　南宋以宋高宗在位第19年作为诏令、赋役与外交文书的纪年框架。此项以《宋史》卷24—47〈本纪〉；《建炎以来系年要录》本年条；《宋会要辑稿》相关门类 为基本来源线索；编年与官修材料能够确认年序和制度用语，但对地方执行、人数、动机或社会后果的判断仍须同文书、碑刻、考古及对方叙事互校。

\paragraph{1146年}\eventanchor{SS-1146-REGNAL-YEAR}{SS-1146-REGNAL-YEAR}　南宋以宋高宗在位第20年作为诏令、赋役与外交文书的纪年框架。此项以《宋史》卷24—47〈本纪〉；《建炎以来系年要录》本年条；《宋会要辑稿》相关门类 为基本来源线索；编年与官修材料能够确认年序和制度用语，但对地方执行、人数、动机或社会后果的判断仍须同文书、碑刻、考古及对方叙事互校。

\paragraph{1147年}\eventanchor{SS-1147-REGNAL-YEAR}{SS-1147-REGNAL-YEAR}　南宋以宋高宗在位第21年作为诏令、赋役与外交文书的纪年框架。此项以《宋史》卷24—47〈本纪〉；《建炎以来系年要录》本年条；《宋会要辑稿》相关门类 为基本来源线索；编年与官修材料能够确认年序和制度用语，但对地方执行、人数、动机或社会后果的判断仍须同文书、碑刻、考古及对方叙事互校。

\paragraph{1148年}\eventanchor{SS-1148-REGNAL-YEAR}{SS-1148-REGNAL-YEAR}　南宋以宋高宗在位第22年作为诏令、赋役与外交文书的纪年框架。此项以《宋史》卷24—47〈本纪〉；《建炎以来系年要录》本年条；《宋会要辑稿》相关门类 为基本来源线索；编年与官修材料能够确认年序和制度用语，但对地方执行、人数、动机或社会后果的判断仍须同文书、碑刻、考古及对方叙事互校。

\paragraph{1149年}\eventanchor{SS-1149-REGNAL-YEAR}{SS-1149-REGNAL-YEAR}　南宋以宋高宗在位第23年作为诏令、赋役与外交文书的纪年框架。此项以《宋史》卷24—47〈本纪〉；《建炎以来系年要录》本年条；《宋会要辑稿》相关门类 为基本来源线索；编年与官修材料能够确认年序和制度用语，但对地方执行、人数、动机或社会后果的判断仍须同文书、碑刻、考古及对方叙事互校。

\paragraph{1150年}\eventanchor{SS-1150-EVENT-122}{SS-1150-EVENT-122}　朝廷调整会子与纸币管理。此项以《宋史》卷24—47〈本纪〉；《建炎以来系年要录》本年条；《宋会要辑稿》相关门类 为基本来源线索；编年与官修材料能够确认年序和制度用语，但对地方执行、人数、动机或社会后果的判断仍须同文书、碑刻、考古及对方叙事互校。

